\documentclass{report}
\usepackage{amsmath,amssymb}
\setlength{\parindent}{0mm}
\setlength{\parskip}{1em}
\begin{document}
\begin{center}
\rule{6in}{1pt} \
{\large David Gardner \\
{\bf Implicit Time Stepping and Preconditioning for Global Atmospheric Dynamics in Climate Simulation}}

Lawrence Livermore National Laboratory \\ 7000 East Ave L-561 \\ Livermore \\ CA 94550
\\
{\tt gardner48@llnl.gov}\\
Richard Archibald\\
Katherine Evans\\
Carol Woodward\end{center}

High resolution, long-term climate simulations are critical to
investigating multiscale behavior in global climate simulations and
understanding regional climate variation on the decade scale. Implicit
time stepping methods provide a means to increase efficiency by taking
time steps commensurate with the physical processes of interest rather
than step sizes limited by model resolution. We present results utilizing
a second order implicit time stepping method and Trilinos solvers in the
hydrostatic spectral element dynamical core of the Community Atmosphere
Model (CAM-SE), the atmosphere component of the Community Earth System
Model (CESM) and the Accelerated Climate Model for Energy (ACME). The
growth of linear iterations in the Newton-Krylov method applied to the
nonlinear residual equation highlights the need for scalable
preconditioning to reduce the cost of ancillary linear system solves. We
discuss the impact of an approximate block factorization preconditioner
on solver efficiency and scalability when utilizing a globally assembled
Jacobian matrix for the spectral element discretization.

This work was performed under the auspices of the U.S. Department of
Energy by Lawrence Livermore National Laboratory under Contract
DE-AC52-07NA27344. LLNL-ABS-681040


\end{document}
